\subsection{命题的逆和否}

\begin{definition}[否命题]
    命题$p$的否为$\neg p$；命题$p\implies q$的否为$\neg p \implies \neg q$。
\end{definition}

\begin{definition}[逆命题]
    命题$p\implies q$的逆为$q \implies p$。
\end{definition}

同时进行逆和否的操作得到的即为\textbf{逆否命题}。

\begin{definition}[逆否命题]
    命题$p\implies q$的逆否为$\neg q \implies \neg p$。-
\end{definition}

\begin{exercise}[逆否命题的真假]
    逆否命题和原命题为等价命题，也即他们同时为真或者同时为假。
\end{exercise}

\begin{solution}
    命题$p\implies q$和$\neg q \implies \neg p$互为逆否命题。

    假设$p\implies q$为真，如果$q$是假命题，那么$p$一定也是假命题，否则通过$p \implies q$可以推出$q$是真命题从而矛盾。所以$\neg q \implies \neg p$。
\end{solution}

\subsection{充分条件和必要条件}
充分条件和必要条件是逻辑蕴涵关系中对命题的叫法。

\begin{definition}[充分条件和必要条件]
    在逻辑蕴含关系$p\implies q$中，$p$是$q$的充分条件，$q$是$p$的必要条件。
\end{definition}

\begin{definition}[充要条件]
    如果$p$既是$q$的充分条件，又是必要条件，那么称为充分必要条件，记作
    $$
        p\iff q
    $$
    也称这两个命题等价。
\end{definition}

\subsection{逻辑量词}

\begin{definition}[存在和任意]
    存在量词“存在”的意思是有一个满足条件即可，符号为$\exists$，全称量词“任意”的意思是对于所有考虑的对象都满足条件，符号为$\forall$。他们都是对命题的修饰，是逻辑量词的一种。
\end{definition}

\begin{exercise}[逻辑量词的否定]
    $$\forall x \in \mathbb{R}, \exists y \in \mathbb{R} \quad x+y=0$$写出这一命题的否定形式。
\end{exercise}

\begin{solution}
    $$\exists x \in \mathbb{R}, \forall y \in \mathbb{R} \quad x+y=0$$

    做否命题的时候，存在量词要改为全称量词，全称量词要改为存在量词。
\end{solution}

\subsection{否命题和命题的否定}

\begin{definition}[命题的否定]
    原命题：$\forall x \leq 0, ~ x \leq 1$
    \begin{enumerate}
        \item 否定：$\exists x \leq 0, x>1$
        \item 否命题：$\forall x>0, x>1$
    \end{enumerate}
\end{definition}

\begin{example}
    设命题 $p: \forall m \in Z, m^2>2 m-3$ ，则 $\neg p$ 为 （\qquad）
    \begin{tasks}(4)
        \task $\forall m \in \mathrm{Z}, m^2 \leq 2 m-3$
        \task $\exists m \in \mathrm{Z}, m^2 \leq 2 m-3$
        \task $\exists m_0 \notin Z, m_0^2>2 m_0-3$
        \task $\forall m \notin \mathrm{Z}, m^2 \leq 2 m-3$
    \end{tasks}
\end{example}

\begin{example}
    命题＂$\forall x \in \mathrm{R}, \exists n \in \mathrm{~N}^*$ ，使得 $n \geq x^2$＂的否定形式是 （\qquad）
    \begin{tasks}(2)
        \task $\forall x \in \mathrm{R}, \exists n \in \mathrm{~N}^*$ ，使得 $n<x^2$
        \task $\forall x \in \mathrm{R}, \forall n \in \mathrm{~N}^*$ ，使得 $n<x^2$
        \task $\exists x \in \mathrm{R}, \exists n \in \mathrm{~N}^*$ ，使得 $n<x^2$
        \task $\exists x_0 \in \mathrm{R}, \forall n \in \mathrm{~N}^*$ ，使得 $n<x_0^2$
    \end{tasks}
\end{example}

